\documentclass[12pt,a4paper]{article}
\usepackage[latin1]{inputenc}
\usepackage[spanish]{babel}
\usepackage{graphicx,amssymb,latexsym,amsmath,verbatim}
\usepackage{sans}
\usepackage{graphicx,amssymb,latexsym,amsmath}
\renewcommand{\thepage}{}
\renewcommand{\baselinestretch}{1}
\setlength{\parindent}{2em} \setlength{\textwidth}{18cm}
\setlength{\textheight}{24cm} \setlength{\topmargin}{-1cm}
\setlength{\oddsidemargin}{-1cm}
%\hyphenation{va-rie-da-des
%va-rie-dad pa-ra-le-lo a-si-mi-la a-pro-ve-cha-bles u-ni-for-me
%u-ni-for-me-men-te si-guien-te si-guien-tes re-su-mien-do
%ha-bi-tua-les}
%%%%%%%%%%%%%%%%%%%%%%%%%%%%%%%%%%%%%%%%%%%%%%%%%
%% Para empezar una interrogativa: ?`  %%%%%%%%%%
%%%%%%%%%%%%%%%%%%%%%%%%%%%%%%%%%%%%%%%%%%%%%%%%%
%% Para partir una matriz en cajas:         %%%%%
%% \[ \left(\begin{array}{c|c} A&B          %%%%%
%% \\ \cline{1-2} C&D\end{array}\right) \]  %%%%%
%%%%%%%%%%%%%%%%%%%%%%%%%%%%%%%%%%%%%%%%%%%%%%%%%


%SETS OF NUMBERS
%\newcommand{\R}{\mbox{${\bf R}$}}                  %reals
\newcommand{\R}{\mbox{${\mathbb R}$}}               %reals
%\newcommand{\C}{\mbox{${\bf C}$}}                  %complex
\newcommand{\C}{\mbox{${\mathbb C}$}}
%\newcommand{\N}{\mbox{${\bf N}$}}                  %natural numbers
\newcommand{\N}{\mbox{${\mathbb N}$}}
%\newcommand{\Z}{\mbox{${\bf Z}$}}                  %integers
\newcommand{\Z}{\mbox{${\mathbb Z}$}}
%\newcommand{\Q}{\mbox{${\bf Q}$}}                  %rationals
\newcommand{\Q}{\mbox{${\mathbb Q}$}}
\newcommand{\gp}{\mbox{${\bf F}_{p}$}}             % Fp
\newcommand{\pea}{\mbox{$ {\bf Z}_{p}$}}           %p-adic integers
\newcommand{\zn}{\mbox{${\bf Z}/{n}$}}             %Z/n
\newcommand{\zp}{\mbox{${\bf Z}/{p}$}}             %Z/p
\newcommand{\Zh}{\mbox{$\hat{{\bf Z}}$}}           %Z^
\newcommand{\xp}{\mbox{${\bf F}_{p}^{*}$}}         %units Fp

%GREEK LETTERS
\newcommand{{\ff}}{\mbox{$\varphi$}}                   %phi
\newcommand{\fp}{\mbox{$\varphi_{p}$}}               %phi_p
\newcommand{\al}{\mbox{$\alpha$}}                    %alpha
\newcommand{\be}{\mbox{$\beta$}}                     %beta
\newcommand{\px}{\mbox{$\varepsilon$}}               %epsilon
\newcommand{\Ga}{\mbox{$\Gamma$}}                    %capital gamma
\newcommand{\ga}{\mbox{$\gamma$}}                    %small gamma

%LOGIC SYMBOLS
\newcommand{\ex}{\mbox{$\exists$}}                   %there is
\newcommand{\fa}{\mbox{$\forall$}}                   %forall
\newcommand{\stf}{\mbox{$\models$}}                  %satisfies
\newcommand{\y}{\mbox{$\wedge$}}                     %and
%\newcommand{\o}{\mbox{$\vee$}}                       %or
\newcommand{\si}{\mbox{$\Longrightarrow$}}           %implies
\newcommand{\ekiv}{\mbox{$\Longleftrightarrow$}}     %equiv

%OTHERS
\newcommand{\dm}{\mbox{$\preceq$}}                   %dominates




\begin{document}

\small

\noindent{\sf Universidad Aut\'onoma de Madrid  \hfill \'Algebra.
Curso 2013--14} \vspace{0.1cm} \hrule \vspace{0.1cm} \noindent{\sf
Ingenier\'ia Inform\'atica} \vspace{0.1cm}
\begin{center}
{\bfseries  \'ALGEBRA. HOJA 5.}
\end{center}

\begin{enumerate}
%\renewcommand{\theenumi}{\bf\arabic{enumi})\ }
\renewcommand{\labelenumi}{\bf\arabic{enumi})\ }
%\renewcommand{\theenumii}{\alph{enumii})\ }
\renewcommand{\labelenumii}{\alph{enumii})\ }
\medskip



\item Aplica el m\'etodo de Gauss para discutir y resolver los siguientes sistemas de ecuaciones lineales.
{\setlength\arraycolsep{.1em}
\begin{align*}
& a)\
\left.
\begin{array}{rrrrrrr}
 && x_2 & -&3x_3 &= & -5\\
2x_1 &+& 3x_2 &+& 3x_3 &=& 7\\
4x_1 &+& 5x_2 &-& 2x_3 &=& 10
\end{array}
\right\}
&&
b)\
\left.
\begin{array}{rrrrrrr}
 3x_1 &-& 10x_2 &-& x_3 &=& -15\\
2x_1  &+& 2x_2 &+& 3x_3 &=& 6\\
 x_1  &+& 14x_2 &+& 7x_3 &=& -1
\end{array}
\right\}
&&
c)\
\left.
\begin{array}{rrrrrrr}
x_1 &+&  x_2 &+& x_3 &=& 2\\
-x_1 &+& x_2 &+& x_3 &=& 1\\
-x_1 &+& 3x_2 &+& 3x_3 &=& 4
\end{array}
\right\}
\\ &
d)\
\left.
\begin{array}{rrrrrrr}
2x_1 &-& 2 x_2 &+&x_3 & = & 9\\
3x_1  &-&5x_2 &+& 2x_3 & = & 4\\
3x_1  &+&3x_2 &-&x_3 & = & 9\\
\end{array}
\right\}
&&
e)\ \left.
\begin{array}{rrrrrrr}
x_1 &-& 2x_2 &+& x_3 &=&7\\
3x_1 &+& 2x_2 &-& x_3 &=& 1\\
2x_1 &-& 5x_2 &+& 2x_3 &=& 6
\end{array}
\right\}
&& f)\ \left.
\begin{array}{rrrrrrr}
x_1 &-& 3x_2 &+& 2x_3  & =&0\\
-x_1 &-& 2x_2 &+& 2x_3  & =& 0\\
2x_1 &&&-& 2x_3  & =& 0
\end{array}
\right\}
\\ &
g)\ \left.
\begin{array}{rrrrrrrrr}
x &+& y &+& z &+& t &=& 0\\
  & & y &-& z & &   &=& 5\\
x &+&   & & z &+& 2t &=& 1\\
x &+& 2y &&   & &   &=& 0
\end{array}
\right\}
&&
h)\ \left.
\begin{array}{rrrrrrr}
x_1 &+& 2x_2 &+& 3x_3 &=& 2\\
x_1 &-& x_2 &+& x_3  &=& 0\\
x_1 &+& 3x_2 &-& x_3 &=& -2\\
3x_1 &+& 4x_2 &+& 3x_3 &=& 0
\end{array}
\right\}
&&
i) \ \left.
\begin{array}{rrrrrrrrr}
x_1 &+& 2x_2 &+& 3x_3 &+& 4x_4 &=& 0\\
 -x_1 &+& x_2 & & &+& 9x_4 &=& 0\\
-x_1 &-& 3x_2 &+& 2 x_3 &+& 3x_4 &=& 0\\
-x_1 &+& 2x_2 &-& 5x_3 &+& 2x_4 &=& 0
\end{array}
\right\}
\end{align*}
}

\noindent
\medskip\textbf{Soluci\'on: } a) $\left( \frac{67}{11}, -\frac{28}{11}, \frac{9}{11}\right)$;
b) incompatible;
c) $\left\{ \left( \frac{1}{2}, \frac{3}{2} - \alpha, \alpha \right) | \alpha \in \R \right\}$;
d) $\left( 1, 13, 33 \right)$;
e) $\left( 2, 8, 21\right)$;\\
f) $\left( 0,0,0\right)$;
g) $\left( -8, 4, -1, 5 \right)$;
h) $\left( -1, 0, 1 \right)$;
i) $\left\{ \left( 59 \alpha, -22 \alpha, -17 \alpha, 9 \alpha \right) | \alpha \in \R \right\}$.

\medskip


\item Discute los siguientes sistemas de ecuaciones en funci\'on de los valores del par\'ametro $a\in\R$:
{\setlength\arraycolsep{.1em}
\[
a) \ \left.
\begin{array}{rrrrrrr}
ax &+& y &+& z &=& 1\\
x &+& ay &+& z &=& a\\
x &+& y &+& az &=& a^2
\end{array}
\right\}
\qquad
b) \
\left.
\begin{array}{rrrrrrr}
3x &-& y && &=& ax\\
5x &+& y &+& 2z &=& ay\\
 && 4y &+& 3z &=& az
\end{array}
\right\}
\]
}

\medskip

\item Discute el siguiente sistema de ecuaciones en funci\'on de los valores de los par\'ametros $a,b\in\R$:
{\setlength\arraycolsep{.1em}
\[
\left.
\begin{array}{rrrrrrr}
2x &-& ay &+& bz &=& 4\\
x  & &    &+& z  &=& 2\\
x &+& y &+& z &=& 2
\end{array}
\right\}
\]
}

\item Cada uno de los siguientes sistemas de ecuaciones tiene dos t\'erminos independientes. Resu\'elvelos a la vez mediante el m\'etodo de Gauss.
{\setlength\arraycolsep{.1em}
\begin{align*}
& a) \ \left.
\begin{array}{rrrrrrrrr|rr}
2x_1 &-& 4x_2 & &     & &     &=& 10 & -8\\
x_1  &-& 3x_2 & &     &+& x_4 &=& -4 & -2\\
x_1  & &      &-& x_3 &+& 2x_4 &=& 4 & 9\\
3x_1 &-& 4x_2 &+& 3x_3 &-& x_4 &=& -11 & -15
\end{array}
\right\} &
b) \ \left.
\begin{array}{rrrrrrr|rr}
2x_1 &-& 4x_2 & & &=& 10 & -8\\
x_1 &-& 3x_2  & & &=& -4 & -2\\
x_1 & & &-& x_3  &=& 4 & 9\\
4x_1 &-& 7x_2 &-& x_3 &=& 10 & -15
\end{array}
\right\}
\\ &
c) \ \left.
\begin{array}{rrrrrrrrr|rr}
2x_1 &-& 4x_2 & & && &=&10 & -8\\
x_1 &-& 3x_2 && &+&x_4 &=& -4 & -2\\
x_1 && &-& x_3 &+& 2x_4 &=& 4 & 9
\end{array}
\right\}
\end{align*}
}

\noindent
\medskip\textbf{Soluci\'on: } a) $\left( \frac{97}{13}, \frac{16}{13}, -\frac{157}{13}, -\frac{101}{13} \right)$ y $\left( 0, 2, -1, 4\right)$;
b) $\left( 23, 9, 19 \right)$ e incompatible;\\
c) $\left\{ \left( 23+2 \alpha, 9+\alpha, 19+4\alpha, \alpha \right) | \alpha \in \R \right\}$ y $\left\{ \left( -8+2\alpha , -2+\alpha, -17+4 \alpha , \alpha \right) | \alpha \in \R \right\}$.

\medskip

\item Aplica el m\'etodo de Gauss para calcular, si existe, la inversa de las siguientes matrices.
\[
a) \ \begin{pmatrix}
1 & 3 & -2\\
2 & 5 & -3\\
-3 & 2 & -4
\end{pmatrix} ,
 \qquad b) \ \begin{pmatrix}
1 & 3 & -2\\
2 & 5 & -3\\
3 & 8 & -5
\end{pmatrix} .
\]
\textbf{Soluci\'on:}
\[
 a) \
 \begin{pmatrix}
14 & -8 & -1\\
-17 & 10 & 1\\
-19 & 11 & 1
\end{pmatrix} , \qquad
b) \ \text{no existe.}
\]

\item Sea $A\in \mathcal{M}_n$ una matriz invertible. Demuestra las siguientes propiedades:

a) $(A^{-1})^{-1}=A$.

b) $(rA)^{-1}=\frac{1}{r}A^{-1}$ para cada $r\in\R \setminus \{0\}$.

c) $(A^p)^{-1}=(A^{-1})^{p}$, para cada $p\in\N$ con $p\geq 1$.

\medskip

\item Dadas dos matrices $A,B\in \mathcal{M}_n$, decide si se cumplen las siguientes igualdades (da una demostraci\'on en caso afirmativo y busca un contraejemplo en caso contrario):

a) $(A+B)^2=A^2+B^2+2AB.$

b) $(A+B)(A-B)=A^2-B^2$.

\medskip

\item Calcula $A^2$, $A^3$, $A^4$ para
\[
 a) \ A = \begin{pmatrix}
1 & 1\\
1 & 1
\end{pmatrix} , \qquad
 b) \ A=
\begin{pmatrix}
1 & 1 & 1\\
0 & 1 & 1\\
0 & 0 & 1
\end{pmatrix} .
\]
Generaliza el resultado calculando $A^n$ para $n\in\N$. (\emph{Sugerencia}: intenta inferir el posible valor de $A^n$ y demu\'estralo por inducci\'on).

\medskip

\item Sea $A=
\begin{pmatrix}
2 & 1\\
1 & 2
\end{pmatrix}$.
Calcula los valores de $\alpha, \beta\in\R$ para los que se cumple la igualdad $A^2+\alpha A+\beta I_2=0$, donde $I_2$ es la matriz identidad de orden $2$.

\medskip

\item Demuestra que la suma de matrices sim\'etricas es sim\'etrica. >Es el producto de matrices sim\'etricas una matriz sim\'etrica en general?


\medskip

\item Una matriz $A\in \mathcal{M}_n$ se llama \emph{idempotente} si $A^2=A$. Demuestra lo siguiente:

a) La matriz $\begin{pmatrix}
2 & -2 & -4\\
-1 & 3 & 4\\
1 & -2 & -3
\end{pmatrix}$ es idempotente.

b) Si $A$ es idempotente, entonces $I_n-A$ es idempotente ($I_n$ es la matriz identidad de orden $n$).

c) Si $A$ es idempotente, entonces $(I_n-A) A=A (I_n-A)=0$.

d) Si $A$ es idempotente e invertible entonces $A$ es la matriz identidad.

\end{enumerate}




\end{document}
